% Generated draft: requires mathematical and source review.
\documentclass[11pt]{article}
\usepackage[margin=25mm]{geometry}
\usepackage{fontspec}
\setmainfont{Latin Modern Roman}
\usepackage{amsmath,amssymb,mathtools,unicode-math}
\setmathfont{Latin Modern Math}
\usepackage{enumitem,needspace,xurl,hyperref}
\setlength{\parindent}{0pt}
\setlength{\parskip}{5pt}
\setlength{\emergencystretch}{3em}
\newcommand{\N}{\mathbb{N}}
\newcommand{\Z}{\mathbb{Z}}
\newcommand{\Q}{\mathbb{Q}}
\newcommand{\R}{\mathbb{R}}
\newcommand{\C}{\mathbb{C}}
\newcommand{\F}{\mathbb{F}}
\newcommand{\A}{\mathbb{A}}
\newcommand{\PP}{\mathbb{P}}
\newcommand{\E}{\mathbb{E}}
\newcommand{\eps}{\varepsilon}
\newcommand{\dd}{\,\mathrm{d}}
\begin{document}
\section*{UnsolvedMath ID: 11000028}
\textbf{Source problem code:} catalogue entry
\subsection{Definitions and mathematical statement}
Let $\operatorname{Mod}_g$ be the mapping class group of a closed orientable genus-$g$ surface. For a finitely generated subgroup $H$ of a finitely generated group $G$, choose finite generating sets $T$ and $S$ respectively and define its distortion function by
\[
 \Delta_H^G(n)=\max\{|h|_T:h\in H,\ |h|_S\le n\}.
\]
Functions are considered up to the standard coarse equivalence: $f\preceq g$ if $f(n)\le Cg(Cn+C)+Cn+C$ for some $C\ge1$, and $f\simeq g$ if both comparisons hold. Linear distortion means $\Delta\simeq n$; precise polynomial distortion of degree $d>1$ means $\Delta\simeq n^d$.

Let $\mathcal I_g$ be the kernel of the action on integral first homology. The source's notation $\mathcal I_g(k)$ appears in a higher-Johnson context. In the customary Johnson-filtration convention, for $\Pi=\pi_1(S_g)$, $\Pi_1=\Pi$, and $\Pi_{j+1}=[\Pi,\Pi_j]$, one has
\[
 \mathcal I_g(k)=\ker\bigl(\operatorname{Mod}_g\to\operatorname{Out}(\Pi/\Pi_{k+1})\bigr),\qquad \mathcal I_g(1)=\mathcal I_g.
\]
The record does not explicitly define this filtration, so this interpretation requires confirmation.

\textbf{Problem 3.7 (distortion).}
\begin{enumerate}
\item Determine the possible distortion types of subgroups of $\operatorname{Mod}_g$.
\item Compute the distortion of $\mathcal I_g$ in $\operatorname{Mod}_g$.
\item Determine how the distortions of $\mathcal I_g(k)$ in $\operatorname{Mod}_g$ behave as $k\to\infty$. The source suggests first computing the successive relative distortions $\Delta_{\mathcal I_g(k+1)}^{\mathcal I_g(k)}$ and suggests higher Johnson homomorphisms as a possible tool.
\item Does there exist a subgroup $H<\operatorname{Mod}_g$ with $\Delta_H^{\operatorname{Mod}_g}(n)\simeq n^d$ for some polynomial degree $d>1$?
\end{enumerate}
The standard definition requires finite generation of each group whose word metric is used. The excerpt gives no genus range or finite-generation convention for all filtration terms. Without those data, the questions concerning them are conditional or require a different specified metric. Also, ``as $k\to\infty$'' refers to a family of functions of $n$; it does not specify uniform estimates or an order of limits. The source guesses exponential distortion for $\mathcal I_g$; that historical guess is not a solution asserted here.

The source reports linear distortion for convex cocompact subgroups, including many free groups [FMo], abelian subgroups [FLMi], and subsurface mapping class subgroups [MM2, Ham]. It also recalls a stronger path condition: for a finite symmetric generating set $S$ of $G$, let $\pi:S^*\to G$ evaluate words and let $\sigma$ choose one or more representative words for each element. A subgroup $H$ is $\sigma$-quasiconvex when a constant $K$ exists such that every path selected by $\sigma(h)$, for every $h\in H$, stays within distance $K$ of $H$ in the Cayley graph. For a uniformly quasigeodesic choice of paths, the source notes that this implies linear distortion [Fa1]. An arbitrary word section is not automatically quasigeodesic.
\subsection{Short English statement}
Classify subgroup distortion in mapping class groups, compute Torelli distortion, understand its Johnson-filtration dependence, and ask whether a genuinely higher-degree polynomial example exists. The filtration questions need finite-generation and metric conventions.
\subsection{Sources}
\begin{itemize}
\item Local UnsolvedMath registry, permanent ID 11000028, source code catalogue entry; no individual proposer or proposal year supplied.
\item Farb, edited volume \emph{Problems on Mapping Class Groups and Related Topics}, Problem 3.7, PDF page 33; \url{https://www.math.uchicago.edu/~farb/papers/mcgbook.pdf}. The registry records PDF-font-structure extraction and access on 2026-07-29; the inventory distinguishes 329 labeled headings by chapter and PDF page.
\item Source references [FMo], [FLMi], [MM2], [Ham], [Mo3], and [Fa1], respectively associated in the excerpt with convex cocompactness, abelian subgroups, subsurface groups, higher Johnson homomorphisms, and quasiconvexity/distortion. Full bibliography entries are not supplied.
\end{itemize}
\end{document}
